\documentclass[11pt, a4paper, oneside]{article}
\usepackage{worksheet}
\usepackage{subfiles}
\usepackage{physics}
\usepackage{siunitx}
\usepackage{comment}

\author{ID:}
\date{17-20.03.2026}
\title{Critical Points worksheet}
\subject{Name: }
\class{Calculus I}

\begin{document}
    \maketitle
    %\testheader[100]{remarks}

    \section*{Instructions}
    The main objective of each of the problems that follow is to determine where \(f'(x) = 0\) or \(f'(x)\) is undefined to identify potential maxima and minima. 
    To determine these points analytically, you will need to use multiple types of differentiation rules, such as the Product, Quotient, and Chain rules. 
    Determine whether the critical parameters for each scenario result in a maximum or minimum.
    Each day, you should finish at least one problem by the end of class.
    The ideal number of completed exercises per session is two, therefore on March 20 you should complete exercises 1-8.
    At the end of the session, please return this page and your process sheets.
    You can utilize your notes to solve the problems.
    Fill in the boxes with your final answers to each task.

	\section*{Exercises}

    \singletask[]{Minimizing Average Cost}
    
    The average of total cost for producing coffee beans is given by 
    \(C(x) = \frac{1}{x}\left[ 200 +10x + \frac{1}{2}x^2 \right]\).
    Compute the minimum average cost of producing coffee beans.

    \boxarea[2cm]

%	\begin{warning}{Warning}
%	Don't be like that	
%	\end{warning}
	
    \singletask[]{Maximum Height of a Projectile}
    
    The height of a proyectile of an object is given by 
    \(h(t) = v_o t -\frac{9.81}{2}t^2\).
    Determine what the object's starting speed was so that its highest point reached at $t=3$.
    Finally, find the maximum height of the proyectile.

    \boxarea[2cm]

    \singletask[]{Volume maximization}
    
    An open box is made from a $3\times 3$ cm sheet by cutting squares of side $x$ from the corner.
    The volume of the box is given by 
    \(V(x) = x(3-2x)^2\).
    Find the length of the squares that achieves the greater volume.

    \boxarea[2cm]

    \singletask[]{Minimizing Surface Area of a Can}
    
    A cylindrical can holds a volume of  \(355~\mathrm{cm}^3\).
    The surface area of this can is given by 
    \(S(r) = 2\pi r^2 + 2V\frac{1}{r}\).
    Find the value of $r$ such that it minimizes the surface area of the can.

    \boxarea[2cm]

    \singletask[]{Finding extremes of Trigonomeric functions}
    
    The displacement of the planets over time can be modeled with the following function,
    \(s(t) = A\sin(\omega t) + B\cos(\omega t) \).
    Find the time at which the planets are in a  maximum or minimum displacements in terms of $A,B$ and $\omega$.

    \boxarea[2cm]

\newpage

    \singletask[]{Maximizing light intensity}

    The intensity of a lamp over a table is given by,
    \( I(\theta) = \sin(\theta)\cos^2(\theta) \).
    Find the critical angle $\theta$ at which the intensity is maximized.

    \boxarea[2cm]

    \grouptask[]{Maximizing Cross-Sectional Area of a Gutter}

    A gutter is bent from a long strip of width $L$ by folding up edges at angle $\theta$.
    The cross-section area is given by \(A(\theta) = (L-2x)x\sin(\theta) + x^2\sin(\theta)\cos(\theta) \).
    Assuming that $L=2$ and $x=1/2$, find the maximum angle $\theta$.

    \boxarea[2cm]

    \singletask[]{Optimal does in medicine}

    The concentration of a drug in blodd after $t$ hours is given by \(C(t) = te^{-kt}\).
    Find the time that gives the maximum concentration in terms of $k$.

    \boxarea[2cm]

\newpage

    \grouptask[]{Catenary}

    A flexible rope of length $L$ suspended between two points at the same height forms a catenary.
    The potential energy of a ball following the rope is given by \(U=a \frac{e^{2x/a}+1}{2e^{x/a}}\).
    Find the minimum of the potential energy.

    \boxarea[2cm]

    \singletask[]{Economic order quantity}

    The total inventory cost per  year is given by \(T(q) = \frac{D}{q}K + \frac{q}{2}h\).
    Where $q$ is the order quantity.
    Find the critical point of the total inventory.

    \boxarea[2cm]

    \singletask[]{Minimizing Fencing Cost}

    A rectangular plot with one side along a river needs fencing on three sides.
    If the area is fixed \(A\) and the cost is given by \(C(x) = px + 2p\frac{A}{x}\), where $x$ is the side along the river.
    Compute side length of the side of the river that minimize the fencing cost given that $A=22500$.

    \boxarea[2cm]

\newpage

    \bonustask[]{Snell's law: Refraction of light}

    A ray of light travels from point $A$ in air (where its speed is $v_1\approx\SI{299 792 458}{\meter\per\second}$) to point $B$ in water (where its speed is $v_1<\SI{299 792 458}{\meter\per\second}$). 
    The interface between the two media is a straight line. 
    The light strikes the interface at a point $x$ units from the point on the shore directly below $A$. 
    If the vertical distance from $A$ to the interface is $a$, and the vertical distance from the interface to $B$ is $b$, the time taken is:
    \begin{gather*}
        T(x) = \frac{\sqrt{x^2+a^2}}{v_1} + \frac{\sqrt{(d-x)^2+b^2}}{v_2}.
    \end{gather*}
    Find \(T'(x)\) and set it to zero to show that the critical point satisfies the condition:
    \begin{gather*}
        \frac{1}{v_1}\sin(\theta_1) = \frac{1}{v_2}\sin(\theta_2),
    \end{gather*}
    where \(\theta_1\) and \(\theta_2\) are the angles of incidence and refraction.
    This is Snell's law.
    Finally show that for $a=b$ and $v_1=2v_3$ the exact critical point is $x=d/3$.

    \section*{Algebra tools and more}

    \begin{minipage}{0.45\textwidth}
        \textbf{Derivatives rules}
    \begin{gather*}
        \dv{x}x^n = nx^{n-1} \\ 
        \dv{x}n^x = n^x\ln(n) \\
        \dv{x}\log_{n}(x) = \frac{1}{x\ln(n)} 
    \end{gather*}
    \end{minipage}
    \hfill
    \begin{minipage}{0.45\textwidth}
        \textbf{Derivatives rules for functions operations}
        \begin{gather*}
            \dv{x}u(x)v(x) = u'(x)v(x) + u(x)v'(x) \\
            \dv{x}\frac{u(x)}{v(x)} = \frac{1}{v^{2}(x)}\left[ u'(x)v(x) - u(x)v'(x) \right] \\
            \dv{x}u(v(x)) = u'(v(x))v'(x) 
        \end{gather*}
    \end{minipage}

    \begin{minipage}{0.45\textwidth}
        \textbf{Properties of exponents}
        \begin{gather*}
            a^{n}\cdot a^{m} = a^{n+m}\\
            \qty(a^{n})^{m} = a^{n\cdot m}\\
            \frac{1}{a} = a^{-1}\\
            a^{\log_{a}\qty[x]} = x = \log_a\qty[a^{x}]
        \end{gather*}
    \end{minipage}
    \hfill
    \begin{minipage}{0.45\textwidth}
        \textbf{Properties of logarithms}
        \begin{gather*}
            \log_{n}\qty[a] + \log_{n}\qty[b] = \log_{n}\qty[a\cdot b] \\
            a\log_{n}\qty[b] = \log_{n}\qty[b^a] \\
            \log_{n}\qty[a] - \log_{n}\qty[b] = \log_{n}\qty[\frac{a}{b}] \\
            \log_{n}\qty[x] = \frac{\log_{m}\qty[x]}{\log_{m}\qty[n]}
        \end{gather*}
    \end{minipage}


\begin{minipage}[t]{0.48\textwidth}
\textbf{Pythagorean Identities}
\begin{align*}
\sin^2\theta + \cos^2\theta &= 1 \\
\tan^2\theta + 1 &= \sec^2\theta \\
1 + \cot^2\theta &= \csc^2\theta
\end{align*}
\vspace{2ex}

\textbf{Sum/Difference Formulas}
\begin{align*}
\sin(\alpha \pm \beta) &= \sin\alpha \cos\beta \pm \cos\alpha \sin\beta \\
\cos(\alpha \pm \beta) &= \cos\alpha \cos\beta \mp \sin\alpha \sin\beta \\
\tan(\alpha \pm \beta) &= \frac{\tan\alpha \pm \tan\beta}{1 \mp \tan\alpha \tan\beta}
\end{align*}
\vspace{2ex}

\textbf{Double-Angle Formulas}
\begin{align*}
\sin(2\theta) &= 2\sin\theta \cos\theta \\
\cos(2\theta) &= \cos^2\theta - \sin^2\theta \\
              &= 2\cos^2\theta - 1 \\
              &= 1 - 2\sin^2\theta \\
\tan(2\theta) &= \frac{2\tan\theta}{1 - \tan^2\theta}
\end{align*}
\vspace{2ex}

\textbf{Periodicity}
\[
\sin(\theta \pm 2\pi) = \sin\theta,\quad 
\cos(\theta \pm 2\pi) = \cos\theta %,\quad 
%\tan(\theta \pm \pi) = \tan\theta
\]
\end{minipage}
\hfill
\begin{minipage}[t]{0.48\textwidth}
\textbf{Half-Angle Formulas}
\begin{align*}
\sin\frac{\theta}{2} &= \pm\sqrt{\frac{1 - \cos\theta}{2}} \\
\cos\frac{\theta}{2} &= \pm\sqrt{\frac{1 + \cos\theta}{2}} \\
\tan\frac{\theta}{2} &= \frac{\sin\theta}{1 + \cos\theta} = \frac{1 - \cos\theta}{\sin\theta}
\end{align*}
\vspace{2ex}

\textbf{Product-to-Sum Formulas}
\begin{align*}
\sin\alpha \sin\beta &= \tfrac{1}{2}[\cos(\alpha - \beta) - \cos(\alpha + \beta)] \\
\cos\alpha \cos\beta &= \tfrac{1}{2}[\cos(\alpha - \beta) + \cos(\alpha + \beta)] \\
\sin\alpha \cos\beta &= \tfrac{1}{2}[\sin(\alpha + \beta) + \sin(\alpha - \beta)]
\end{align*}
\vspace{2ex}

\textbf{Sum-to-Product Formulas}
\begin{align*}
\sin\alpha + \sin\beta &= 2\sin\frac{\alpha + \beta}{2}\cos\frac{\alpha - \beta}{2} \\
\sin\alpha - \sin\beta &= 2\cos\frac{\alpha + \beta}{2}\sin\frac{\alpha - \beta}{2} \\
\cos\alpha + \cos\beta &= 2\cos\frac{\alpha + \beta}{2}\cos\frac{\alpha - \beta}{2} \\
\cos\alpha - \cos\beta &= -2\sin\frac{\alpha + \beta}{2}\sin\frac{\alpha - \beta}{2}
\end{align*}
\vspace{2ex}


\textbf{Cofunction Identities}
\[
\sin\left(\frac{\pi}{2} - \theta\right) = \cos\theta,\quad
\cos\left(\frac{\pi}{2} - \theta\right) = \sin\theta
\]
\end{minipage}


\begin{comment}
\textbf{Reciprocal Identities}
\[
\csc\theta = \frac{1}{\sin\theta},\quad 
\sec\theta = \frac{1}{\cos\theta},\quad 
\cot\theta = \frac{1}{\tan\theta}
\]

\vspace{2ex}



    \subsection*{Trigonometric Identities}
\begin{enumerate}
    \item \textbf{Pythagorean Identities:}
    \begin{align*}
        \sin^2\theta + \cos^2\theta &= 1 \\
        \tan^2\theta + 1 &= \sec^2\theta \\
        1 + \cot^2\theta &= \csc^2\theta
    \end{align*}

    \item \textbf{Sum/Difference Formulas:}
    \begin{align*}
        \sin(\alpha \pm \beta) &= \sin\alpha \cos\beta \pm \cos\alpha \sin\beta \\
        \cos(\alpha \pm \beta) &= \cos\alpha \cos\beta \mp \sin\alpha \sin\beta \\
        \tan(\alpha \pm \beta) &= \frac{\tan\alpha \pm \tan\beta}{1 \mp \tan\alpha \tan\beta}
    \end{align*}

    \item \textbf{Double-Angle Formulas:}
    \begin{align*}
        \sin(2\theta) &= 2\sin\theta \cos\theta \\
        \cos(2\theta) &= \cos^2\theta - \sin^2\theta \\
                      &= 2\cos^2\theta - 1 \\
                      &= 1 - 2\sin^2\theta \\
        \tan(2\theta) &= \frac{2\tan\theta}{1 - \tan^2\theta}
    \end{align*}

    \item \textbf{Half-Angle Formulas:}
    \begin{align*}
        \sin\frac{\theta}{2} &= \pm\sqrt{\frac{1 - \cos\theta}{2}} \\
        \cos\frac{\theta}{2} &= \pm\sqrt{\frac{1 + \cos\theta}{2}} \\
        \tan\frac{\theta}{2} &= \frac{\sin\theta}{1 + \cos\theta} = \frac{1 - \cos\theta}{\sin\theta}
    \end{align*}

    \item \textbf{Product-to-Sum Formulas:}
    \begin{align*}
        \sin\alpha \sin\beta &= \frac{1}{2}[\cos(\alpha - \beta) - \cos(\alpha + \beta)] \\
        \cos\alpha \cos\beta &= \frac{1}{2}[\cos(\alpha - \beta) + \cos(\alpha + \beta)] \\
        \sin\alpha \cos\beta &= \frac{1}{2}[\sin(\alpha + \beta) + \sin(\alpha - \beta)]
    \end{align*}

    \item \textbf{Sum-to-Product Formulas:}
    \begin{align*}
        \sin\alpha + \sin\beta &= 2\sin\frac{\alpha + \beta}{2}\cos\frac{\alpha - \beta}{2} \\
        \sin\alpha - \sin\beta &= 2\cos\frac{\alpha + \beta}{2}\sin\frac{\alpha - \beta}{2} \\
        \cos\alpha + \cos\beta &= 2\cos\frac{\alpha + \beta}{2}\cos\frac{\alpha - \beta}{2} \\
        \cos\alpha - \cos\beta &= -2\sin\frac{\alpha + \beta}{2}\sin\frac{\alpha - \beta}{2}
    \end{align*}

    \item \textbf{Reciprocal Identities:}
    \[
        \csc\theta = \frac{1}{\sin\theta},\quad 
        \sec\theta = \frac{1}{\cos\theta},\quad 
        \cot\theta = \frac{1}{\tan\theta}
    \]

    \item \textbf{Periodicity:}
    \begin{align*}
        \sin(\theta \pm 2\pi) &= \sin\theta, &
        \cos(\theta \pm 2\pi) &= \cos\theta, \\
        \tan(\theta \pm \pi) &= \tan\theta
    \end{align*}

    \item \textbf{Cofunction Identities:}
    \[
        \sin\left(\frac{\pi}{2} - \theta\right) = \cos\theta,\quad
        \cos\left(\frac{\pi}{2} - \theta\right) = \sin\theta
    \]
\end{enumerate}
\end{comment}    

%        \partnertask[10]{This is a title}

%    \grouptask[10]{This is a title}

%    \checkered[0.85\textwidth]

%    \bonustask[10]{This is a title}
    
%    \lines[0.85\textwidth]

%    \testtask[10]{Miau}
    
%    \boxarea[0.85\textwidth]

\end{document}
